%% ========================================================================
%% week04-filesystem-ops.tex
%% Modul Praktikum: Struktur Direktori & Operasi File pada Linux
%% Target: Ubuntu Server 22.04
%% Gaya penulisan mengikuti pola week03-basic-io.tex:
%%   - chapter + abstract
%%   - penjelasan konsep singkat namun jelas
%%   - daftar perintah (fungsi)
%%   - praktikum minimal 2 per topik (mudah + kondisi riil/sulit)
%%   - latihan di akhir modul
%%
%% Catatan:
%% - Environment seperti tipbox, warningbox, notebox, examplebox, exercisebox
%%   serta style lstlisting diasumsikan sudah tersedia dari preamble master,
%%   sama seperti pada week03-basic-io.tex.
%% - Jika environment box belum ada, Anda bisa ganti tipbox/warningbox/...
%%   menjadi \subsubsection*{} biasa.
%% ========================================================================

\chapter{Struktur Direktori dan Operasi File pada Linux}
\label{ch:filesystem-ops}

Bab ini membahas struktur direktori Linux berdasarkan praktik umum \textit{Filesystem Hierarchy Standard} (FHS),
jenis-jenis file di Linux (reguler, direktori, symlink, dan file perangkat), navigasi sistem file melalui command line,
operasi file lanjutan terkait izin dan kepemilikan (\texttt{chmod}, \texttt{chown}, \texttt{chgrp}),
hard links dan symbolic links, pencarian file dengan \texttt{find} dan \texttt{locate},
serta pengarsipan dan kompresi (\texttt{tar}, \texttt{gzip}, \texttt{zip}).
Modul dirancang untuk praktikum Ubuntu Server 22.04 dengan pendekatan \textit{hands-on}: konsep, verifikasi output,
skenario mudah, dan skenario kondisi riil yang sering ditemui administrator server.

% ------------------------------------------------------------------------
\section{Standar Hierarki Sistem File (FHS) di Linux}
\label{sec:fhs}

\subsection*{Penjelasan Konsep}
Linux menyusun file dan folder dalam struktur hirarkis yang berpusat pada root (\texttt{/}).
Pada server, memahami \textit{“apa ada di mana”} adalah kemampuan dasar: konfigurasi ada di \texttt{/etc},
log banyak berada di \texttt{/var/log}, home user ada di \texttt{/home}, dan perangkat direpresentasikan sebagai file di \texttt{/dev}.
Prinsip Unix/Linux yang relevan adalah \textit{everything is a file}: bukan hanya file reguler,
tetapi juga perangkat (\texttt{/dev}) dan informasi sistem/proses (\texttt{/proc}, \texttt{/sys}) diperlakukan seperti file.

\begin{tipbox}
Cara berpikir yang membantu saat administrasi server:
\textit{``konfigurasi itu file'', ``log itu file'', ``device itu file''}.
Maka banyak troubleshooting berakhir pada \texttt{cat}, \texttt{less}, \texttt{grep}, dan \texttt{ls -l}.
\end{tipbox}

\subsection*{Direktori penting (ringkas + fungsi)}
\begin{itemize}
    \item \texttt{/} : root dari seluruh filesystem.
    \item \texttt{/bin} : program esensial untuk semua user (contoh: \texttt{ls}, \texttt{cp}, \texttt{mv}).
    \item \texttt{/sbin} : program esensial administrasi (sering butuh \texttt{sudo}).
    \item \texttt{/etc} : konfigurasi sistem (umumnya file teks).
    \item \texttt{/var} : data yang sering berubah (log, cache, spool). Umum: \texttt{/var/log}.
    \item \texttt{/home} : direktori home user.
    \item \texttt{/root} : home untuk user \texttt{root}.
    \item \texttt{/tmp} : file sementara (sering dibersihkan otomatis).
    \item \texttt{/usr} : aplikasi dan library userland (sub: \texttt{/usr/bin}, \texttt{/usr/sbin}).
    \item \texttt{/opt} : aplikasi pihak ketiga/opsional (sering dipakai untuk deploy aplikasi custom).
    \item \texttt{/srv} : data untuk layanan server (mis. web/app).
    \item \texttt{/dev} : \textit{device node} (disk, tty, dll).
    \item \texttt{/proc}, \texttt{/sys} : \textit{virtual filesystem} dari kernel (status proses/sistem).
    \item \texttt{/mnt}, \texttt{/media} : mount point sementara/removable.
\end{itemize}

\subsection*{Praktikum 4.1 --- Tur FHS: Kenali Isi Sistem (Mudah)}
\textbf{Tujuan:} mengenali folder penting dan membedakan mana konfigurasi, log, dan data user.

\begin{enumerate}
    \item Cek posisi Anda dan tampilkan isi root:
\begin{lstlisting}[language=bash, caption={Melihat root filesystem dan isinya}]
pwd
ls -lah /
\end{lstlisting}

    \item Lihat isi folder konfigurasi dan log (ringkas):
\begin{lstlisting}[language=bash, caption={Menginspeksi /etc dan /var/log}]
ls -lah /etc | head
ls -lah /var/log | head
\end{lstlisting}

    \item (Opsional) Pasang \texttt{tree} untuk visual hirarki:
\begin{lstlisting}[language=bash, caption={Memasang tree dan menampilkan struktur folder}]
sudo apt update
sudo apt install -y tree
tree -L 2 /
\end{lstlisting}
\end{enumerate}

\begin{warningbox}
Jangan menjalankan \texttt{tree /} tanpa batas kedalaman pada server nyata—output bisa sangat besar.
Gunakan \texttt{-L} untuk membatasi kedalaman.
\end{warningbox}

\subsection*{Praktikum 4.2 --- Audit Lokasi Konfigurasi dan Log Layanan (Kondisi Riil)}
\textbf{Skenario:} admin diminta mengecek konfigurasi SSH dan melihat indikasi akses pada log.

\textbf{Tujuan:} membiasakan mencari file konfigurasi standar dan membaca log secara aman.

\begin{enumerate}
    \item Verifikasi file konfigurasi SSH:
\begin{lstlisting}[language=bash, caption={Cek konfigurasi SSH}]
ls -lah /etc/ssh/
sudo test -f /etc/ssh/sshd_config \
  && echo "sshd_config ada" \
  || echo "sshd_config tidak ditemukan"
\end{lstlisting}

    \item Baca baris-baris awal konfigurasi (tanpa mengedit):
\begin{lstlisting}[language=bash, caption={Membaca konfigurasi dengan aman}]
sudo head -n 30 /etc/ssh/sshd_config
\end{lstlisting}

    \item Lihat aktivitas terakhir pada log autentikasi:
\begin{lstlisting}[language=bash, caption={Melihat log autentikasi}]
ls -lah /var/log/ | grep -E "auth|secure" || true
sudo tail -n 50 /var/log/auth.log
\end{lstlisting}

    \item (Opsional) Filter event penting:
\begin{lstlisting}[language=bash, caption={Filter event penting pada auth.log}]
sudo grep \ 
  -E "Failed password|Accepted password|Invalid user" \
  /var/log/auth.log | tail -n 30
\end{lstlisting}
\end{enumerate}

% ------------------------------------------------------------------------
\section{Jenis File di Linux (Reguler, Direktori, Symlink, Perangkat)}
\label{sec:filetypes}

\subsection*{Penjelasan Konsep}
Di Linux, ``file'' bukan cuma dokumen. Dari output \texttt{ls -l}, karakter pertama menunjukkan jenis:
\begin{itemize}
    \item \texttt{-} : file reguler
    \item \texttt{d} : direktori
    \item \texttt{l} : symbolic link
    \item \texttt{b} : block device (disk/partisi)
    \item \texttt{c} : character device (tty/serial)
    \item \texttt{p} : FIFO (named pipe)
    \item \texttt{s} : socket
\end{itemize}

\subsection*{Perintah yang dipakai}
\begin{itemize}
    \item \texttt{ls -l} / \texttt{ls -lah}: melihat jenis file dan permission.
    \item \texttt{stat <file>}: detail metadata (inode, size, waktu, permission).
    \item \texttt{file <file>}: menebak tipe konten (ASCII text, image, ELF binary, dll).
\end{itemize}

\subsection*{Praktikum 4.3 --- Membaca Jenis File dari Output ls -l (Mudah)}
\textbf{Tujuan:} membedakan file reguler, direktori, symlink, dan device.

\begin{enumerate}
    \item Buat workspace praktikum:
\begin{lstlisting}[language=bash, caption={Membuat workspace bab ini}]
mkdir -p ~/praktikum-os/week04
cd ~/praktikum-os/week04
pwd
\end{lstlisting}

    \item Buat contoh file dan direktori:
\begin{lstlisting}[language=bash, caption={Membuat file reguler dan direktori contoh}]
mkdir -p data
echo "hello server" > data/hello.txt
ls -lah
ls -lah data
\end{lstlisting}

    \item Buat symlink sederhana:
\begin{lstlisting}[language=bash, caption={Membuat symlink}]
ln -s data/hello.txt hello-link.txt
ls -lah
\end{lstlisting}

    \item Inspeksi metadata:
\begin{lstlisting}[language=bash, caption={Menginspeksi metadata file}]
stat data/hello.txt | head -n 20
file data/hello.txt
\end{lstlisting}
\end{enumerate}

\subsection*{Praktikum 4.4 --- Membaca File Perangkat dan Virtual FS (Kondisi Riil)}
\textbf{Skenario:} admin perlu membedakan device node (block vs char) dan file virtual dari kernel.

\textbf{Tujuan:} mengenali \texttt{/dev} (device node), \texttt{/proc} dan \texttt{/sys} (virtual filesystem).

\begin{enumerate}
    \item Identifikasi block device dan character device:
\begin{lstlisting}[language=bash, caption={Block vs character device}]
ls -l /dev/sd* 2>/dev/null | head -n 10
ls -l /dev/tty
\end{lstlisting}

    \item Tunjukkan bahwa \texttt{/proc} berisi info proses (virtual):
\begin{lstlisting}[language=bash, caption={Membaca informasi sistem dari /proc}]
ls -lah /proc | head
cat /proc/cpuinfo | head -n 20
\end{lstlisting}

    \item Gunakan \texttt{file} untuk membedakan konten:
\begin{lstlisting}[language=bash, caption={Menguji tipe konten dengan file}]
file /dev/tty
file /proc/cpuinfo
\end{lstlisting}
\end{enumerate}

\begin{warningbox}
Jangan menulis sembarang ke \texttt{/dev/*} atau \texttt{/proc/*}. Banyak entry bersifat sensitif dan bisa memengaruhi sistem.
\end{warningbox}

% ------------------------------------------------------------------------
\section{Navigasi Sistem File melalui Command Line}
\label{sec:navigasi}

\subsection*{Penjelasan Konsep}
Pada Ubuntu Server, terminal adalah antarmuka utama administrasi.
Skill navigasi (menemukan lokasi file, berpindah folder, dan membaca struktur) adalah fondasi semua tugas lanjutan.

\subsection*{Perintah yang dipakai}
\begin{itemize}
    \item \texttt{pwd}: menampilkan direktori aktif.
    \item \texttt{ls}: daftar isi direktori (opsi umum: \texttt{-l}, \texttt{-a}, \texttt{-h}).
    \item \texttt{cd}: pindah direktori.
    \item \texttt{pushd}/\texttt{popd}: navigasi stack direktori (nyaman untuk bolak-balik).
\end{itemize}

\subsection*{Praktikum 4.5 --- Navigasi Aman dan Efisien (Mudah)}
\textbf{Tujuan:} membiasakan absolute path, relative path, dan \texttt{pushd/popd}.

\begin{enumerate}
    \item Pindah ke root, lalu kembali ke home:
\begin{lstlisting}[language=bash, caption={Absolute path vs kembali ke home}]
cd /
pwd
cd ~
pwd
\end{lstlisting}

    \item Latihan relative path:
\begin{lstlisting}[language=bash, caption={Relative path di workspace}]
cd ~/praktikum-os/week04
mkdir -p a/b/c
cd a/b
pwd
cd ../..
pwd
\end{lstlisting}

    \item Gunakan \texttt{pushd/popd} untuk berpindah bolak-balik direktori:
\begin{lstlisting}[language=bash, caption={Navigasi stack direktori}]
pushd /etc
pwd
popd
pwd
\end{lstlisting}
\end{enumerate}

\begin{warningbox}
Perintah \texttt{rm -rf} dapat menghapus direktori beserta isinya tanpa konfirmasi.
Untuk praktikum, biasakan memakai \texttt{rm -ri} atau pastikan path benar.
\end{warningbox}

\subsection*{Praktikum 4.6 --- Menemukan File Konfigurasi dari Output Service (Kondisi Riil)}
\textbf{Skenario:} Anda diminta mencari lokasi unit service dan file konfigurasi terkait.

\textbf{Tujuan:} menggunakan navigasi + pembacaan path sistem yang umum pada server.

\begin{enumerate}
    \item Lihat status SSH dan petunjuk lokasi:
\begin{lstlisting}[language=bash, caption={Menelusuri info layanan via systemctl}]
systemctl status ssh --no-pager | head -n 30
\end{lstlisting}

    \item Tampilkan unit file layanan:
\begin{lstlisting}[language=bash, caption={Menemukan unit file layanan}]
systemctl cat ssh --no-pager | head -n 80
\end{lstlisting}

    \item Navigasi cepat ke direktori konfigurasi:
\begin{lstlisting}[language=bash, caption={Navigasi ke direktori konfigurasi SSH}]
cd /etc/ssh
pwd
ls -lah | head
\end{lstlisting}
\end{enumerate}

% ------------------------------------------------------------------------
\section{Operasi File Lanjutan: chmod, chown, chgrp}
\label{sec:perm-owner}

\subsection*{Penjelasan Konsep}
Kontrol akses di Linux dibangun dari tiga komponen:
\begin{itemize}
    \item \textbf{Permission} (izin) untuk \textit{owner}, \textit{group}, \textit{others}
    \item \textbf{Ownership} (kepemilikan): \textit{user owner} dan \textit{group owner}
    \item \textbf{Konsep eksekusi direktori}: pada direktori, \texttt{x} berarti ``boleh masuk/akses isi direktori'' (traverse).
\end{itemize}

\subsection*{Membaca permission: rwx untuk user/group/others}
Output \texttt{ls -l} contoh:
\begin{center}
\texttt{-rwxr-x---}
\end{center}
Artinya:
\begin{itemize}
    \item Karakter 1: jenis file (\texttt{-}, \texttt{d}, \texttt{l}, dst)
    \item 3 berikutnya: \textbf{owner} (user)
    \item 3 berikutnya: \textbf{group}
    \item 3 terakhir: \textbf{others}
\end{itemize}

\subsubsection*{Makna r/w/x pada file vs direktori}
\begin{itemize}
    \item \textbf{File reguler}:
    \begin{itemize}
        \item \texttt{r}: boleh membaca isi file
        \item \texttt{w}: boleh mengubah isi file
        \item \texttt{x}: boleh mengeksekusi file sebagai program/script
    \end{itemize}
    \item \textbf{Direktori}:
    \begin{itemize}
        \item \texttt{r}: boleh melihat daftar nama file (mis. \texttt{ls})
        \item \texttt{w}: boleh membuat/menghapus/rename entry di direktori
        \item \texttt{x}: boleh masuk (\texttt{cd}) dan mengakses file di dalamnya
    \end{itemize}
\end{itemize}

\begin{warningbox}
Direktori tanpa \texttt{x} bisa terlihat isinya (jika \texttt{r} ada), tetapi tidak bisa diakses/masuk.
Sebaliknya, direktori dengan \texttt{x} tapi tanpa \texttt{r} bisa diakses jika Anda tahu nama file tepatnya.
\end{warningbox}

\subsection{chmod: Notasi Angka (Oktal)}
\label{subsec:chmod-oktal}

Setiap permission direpresentasikan sebagai angka oktal (0--7). Setiap triad
(\textit{user/group/others}) dihitung dengan menjumlahkan nilai bit:

\begin{center}
\begin{tabular}{ccc}
    \textbf{r} (read) & \textbf{w} (write) & \textbf{x} (execute) \\
    4                 & 2                  & 1                    \\
\end{tabular}
\end{center}

Jumlahkan nilai bit yang aktif untuk mendapatkan angka oktal tiap triad.
Tiga angka oktal digabung membentuk permission lengkap, misalnya \texttt{755}
berarti user=7, group=5, others=5.

\begin{table}[htbp]
    \centering
    \caption{Tabel konversi notasi oktal ke permission}
    \begin{tabular}{cll}
        \hline
        \textbf{Angka} & \textbf{Izin} & \textbf{Deskripsi} \\
        \hline
        0 & \texttt{-{}-{}-} & Tidak ada akses sama sekali \\
        1 & \texttt{-{}-x} & Hanya eksekusi (masuk direktori) \\
        2 & \texttt{-w-} & Hanya tulis \\
        3 & \texttt{-wx} & Tulis dan eksekusi \\
        4 & \texttt{r-{}-} & Hanya baca \\
        5 & \texttt{r-x} & Baca dan eksekusi \\
        6 & \texttt{rw-} & Baca dan tulis \\
        7 & \texttt{rwx} & Baca, tulis, dan eksekusi (akses penuh) \\
        \hline
    \end{tabular}
\end{table}

Contoh penggunaan:
\begin{itemize}
    \item \texttt{chmod 755 file} $\rightarrow$ owner \texttt{rwx} (7), group \texttt{r-x} (5), others \texttt{r-x} (5)
    \item \texttt{chmod 640 file} $\rightarrow$ owner \texttt{rw-} (6), group \texttt{r--} (4), others \texttt{---} (0)
    \item \texttt{chmod 600 file} $\rightarrow$ owner \texttt{rw-} (6), group \texttt{---} (0), others \texttt{---} (0)
\end{itemize}

\subsection*{chmod: notasi simbolik}
\begin{itemize}
    \item Target: \texttt{u} (user), \texttt{g} (group), \texttt{o} (others), \texttt{a} (all)
    \item Operator: \texttt{+} (tambah), \texttt{-} (hapus), \texttt{=} (set persis)
\end{itemize}

Contoh:
\begin{itemize}
    \item \texttt{chmod u+x script.sh} (tambah execute untuk owner)
    \item \texttt{chmod go-rw file} (hapus read+write untuk group dan others)
\end{itemize}

\subsection*{Special bits: setuid, setgid, sticky}
Selain 3 triad utama, ada bit khusus:
\begin{itemize}
    \item \textbf{setuid (4xxx)}: program berjalan dengan hak owner file (umumnya root). Contoh: \texttt{/usr/bin/passwd}.
    \item \textbf{setgid (2xxx)}:
    \begin{itemize}
        \item pada file: berjalan dengan group file
        \item pada direktori: file baru mewarisi group direktori (\textit{penting untuk shared folder})
    \end{itemize}
    \item \textbf{sticky bit (1xxx)} pada direktori: hanya owner file (atau root) yang boleh menghapus file di direktori tersebut.
          Contoh: \texttt{/tmp}.
\end{itemize}

\subsection*{Ownership: chown dan chgrp}
\begin{itemize}
    \item \texttt{chown user file} mengubah owner user
    \item \texttt{chown user:group file} mengubah user dan group sekaligus
    \item \texttt{chgrp group file} mengubah group
\end{itemize}

\begin{tipbox}
Masalah akses pada server sering bukan karena file ``hilang'', tetapi karena ownership/permission tidak sesuai.
Biasakan verifikasi dengan \texttt{ls -l} dan \texttt{stat}.
\end{tipbox}

\subsection*{Praktikum 4.7 --- Permission Dasar dengan File dan Direktori (Mudah)}
\textbf{Tujuan:} melihat perbedaan makna \texttt{x} pada file vs direktori, serta latihan konversi oktal.

\begin{enumerate}
    \item Siapkan objek uji:
\begin{lstlisting}[language=bash, caption={Setup objek uji permission}]
cd ~/praktikum-os/week04
mkdir -p permtest/dirA
echo "hello" > permtest/fileA.txt
ls -ld permtest permtest/dirA
ls -l permtest/fileA.txt
\end{lstlisting}

    \item Ubah permission file menjadi \texttt{640} lalu verifikasi:
\begin{lstlisting}[language=bash, caption={chmod file dengan oktal}]
chmod 640 permtest/fileA.txt
ls -l permtest/fileA.txt
\end{lstlisting}

    \item Ubah permission direktori menjadi \texttt{700}, lalu coba akses:
\begin{lstlisting}[language=bash, caption={chmod direktori dan uji akses}]
chmod 700 permtest/dirA
ls -ld permtest/dirA
cd permtest/dirA
pwd
cd -
\end{lstlisting}

    \item Eksperimen: hilangkan \texttt{x} dari direktori dan lihat efeknya:
\begin{lstlisting}[language=bash, caption={Menguji efek hilangnya x pada direktori}]
chmod 600 permtest/dirA
ls -ld permtest/dirA
cd permtest/dirA 2>/dev/null \
 || echo "Gagal: tidak ada execute (x) pada direktori"
\end{lstlisting}
\end{enumerate}

\subsection*{Praktikum 4.8 --- Shared Folder Proyek + Sticky Bit (Kondisi Riil)}
\textbf{Skenario:} folder proyek bersama: anggota group bisa kolaborasi, file baru mewarisi group,
dan tidak saling menghapus file sembarangan.

\textbf{Tujuan:} menerapkan \texttt{chgrp}, \texttt{chmod 2xxx}, dan sticky bit pada direktori kolaboratif.

\begin{enumerate}
    \item Buat group dan folder proyek:
\begin{lstlisting}[language=bash, caption={Membuat group dan folder proyek}]
sudo groupadd project 2>/dev/null || echo "group project sudah ada"
sudo usermod -aG project $USER
sudo mkdir -p /srv/project-gamma
sudo chgrp project /srv/project-gamma
\end{lstlisting}

    \item Terapkan setgid agar file baru mewarisi group, dan batasi others:
\begin{lstlisting}[language=bash, caption={Setgid pada direktori proyek}]
sudo chmod 2770 /srv/project-gamma
ls -ld /srv/project-gamma
\end{lstlisting}

    \item Terapkan sticky bit agar user tidak bisa menghapus file milik user lain:
\begin{lstlisting}[language=bash, caption={Sticky bit pada direktori kolaboratif}]
sudo chmod +t /srv/project-gamma
ls -ld /srv/project-gamma
\end{lstlisting}

    \item Uji pembuatan file (catatan: membership group bisa butuh re-login):
\begin{lstlisting}[language=bash, caption={Uji akses kolaborasi}]
touch /srv/project-gamma/test-$USER.txt 2>/dev/null \
  || echo "Jika gagal: re-login atau jalankan newgrp"
ls -l /srv/project-gamma | head
\end{lstlisting}
\end{enumerate}

\begin{exercisebox}[4.1]
Tentukan angka chmod untuk:
\begin{enumerate}
    \item File: owner \texttt{rw-}, group \texttt{r--}, others \texttt{---}
    \item Direktori shared: owner+group \texttt{rwx}, others \texttt{---}, file baru mewarisi group (setgid)
    \item Direktori seperti \texttt{/tmp}: semua bisa buat file, tetapi tidak bisa hapus file user lain (sticky bit)
\end{enumerate}
\end{exercisebox}

% ------------------------------------------------------------------------
\section{Hard Links dan Symbolic Links}
\label{sec:links}

\subsection*{Penjelasan Konsep}
Link adalah cara membuat ``nama lain'' untuk sebuah file.
\begin{itemize}
    \item \textbf{Hard link} (\texttt{ln}): menunjuk ke inode yang sama. Menghapus nama asli tidak menghapus data jika masih ada hard link lain.
    \item \textbf{Symbolic link} / symlink (\texttt{ln -s}): file kecil yang menunjuk ke path target. Jika target pindah/hilang, symlink menjadi \textit{broken}.
\end{itemize}

\subsection*{Perintah yang dipakai}
\begin{itemize}
    \item \texttt{ln}, \texttt{ln -s}: membuat link.
    \item \texttt{ls -li}: menampilkan inode dan link count.
    \item \texttt{stat}: verifikasi inode dan metadata.
\end{itemize}

\subsection*{Praktikum 4.9 --- Uji Hard Link vs Symlink (Mudah)}
\textbf{Tujuan:} melihat perbedaan inode dan efek penghapusan.

\begin{enumerate}
    \item Siapkan file:
\begin{lstlisting}[language=bash, caption={Membuat file sumber}]
cd ~/praktikum-os/week04
echo "version1" > target.txt
\end{lstlisting}

    \item Buat hard link dan symlink:
\begin{lstlisting}[language=bash, caption={Membuat hard link dan symlink}]
ln target.txt hardlink.txt
ln -s target.txt symlink.txt
ls -li target.txt hardlink.txt symlink.txt
\end{lstlisting}

    \item Hapus file asli dan cek efeknya:
\begin{lstlisting}[language=bash, caption={Menghapus target dan mengecek link}]
rm target.txt
ls -li hardlink.txt symlink.txt
cat hardlink.txt
cat symlink.txt \
  || echo "symlink broken karena target hilang"
\end{lstlisting}
\end{enumerate}

\begin{warningbox}
Hard link umumnya tidak dipakai untuk direktori (dibatasi untuk mencegah loop struktur).
Untuk kebutuhan ``shortcut'' direktori, gunakan symlink.
\end{warningbox}

\subsection*{Praktikum 4.10 --- Symlink untuk Deploy Versi Aplikasi (Kondisi Riil)}
\textbf{Skenario:} pola deployment umum: folder versi aplikasi terpisah, lalu \texttt{current} adalah symlink ke versi aktif.

\textbf{Tujuan:} memahami kenapa symlink sangat berguna untuk rollback cepat.

\begin{enumerate}
    \item Buat struktur versi:
\begin{lstlisting}[language=bash, caption={Membuat versi aplikasi palsu}]
cd ~/praktikum-os/week04
mkdir -p app/releases/v1 app/releases/v2 app/shared
echo "APP_VERSION=v1" > app/releases/v1/.env
echo "APP_VERSION=v2" > app/releases/v2/.env
\end{lstlisting}

    \item Buat symlink \texttt{current} menunjuk ke versi aktif:
\begin{lstlisting}[language=bash, caption={Membuat symlink current}]
ln -s releases/v1 app/current
ls -lah app
cat app/current/.env
\end{lstlisting}

    \item Simulasikan upgrade dan rollback hanya dengan mengganti symlink:
\begin{lstlisting}[language=bash, caption={Switch versi via symlink}]
rm app/current
ln -s releases/v2 app/current
cat app/current/.env

rm app/current
ln -s releases/v1 app/current
cat app/current/.env
\end{lstlisting}
\end{enumerate}

% ------------------------------------------------------------------------
\section{Pencarian File dengan find dan locate}
\label{sec:search}

\subsection*{Penjelasan Konsep}
Pada server, Anda sering perlu mencari:
\begin{itemize}
    \item file konfigurasi (\texttt{*.conf}), log, atau file besar yang memakan disk
    \item file yang permission-nya bermasalah
    \item file yang berubah dalam rentang waktu tertentu (debug)
\end{itemize}

\subsection*{Perintah yang dipakai}
\begin{itemize}
    \item \texttt{find}: pencarian real-time (akurat, fleksibel).
    \item \texttt{locate}: pencarian berbasis database (cepat, perlu update database).
\end{itemize}

\subsection*{Praktikum 4.11 --- find Dasar: Nama, Tipe, Ukuran (Mudah)}
\textbf{Tujuan:} mencari file berdasarkan nama, tipe, dan ukuran.

\begin{enumerate}
    \item Cari semua file \texttt{.txt} di workspace:
\begin{lstlisting}[language=bash, caption={find berdasarkan nama dan tipe}]
cd ~/praktikum-os/week04
find . -type f -name "*.txt"
\end{lstlisting}

    \item Cari file lebih besar dari 1MB (contoh umum troubleshooting):
\begin{lstlisting}[language=bash, caption={find berdasarkan ukuran file}]
find ~ -type f -size +1M 2>/dev/null | head
\end{lstlisting}
\end{enumerate}

\subsection*{Praktikum 4.12 --- find Kondisi Riil: Scan Konfigurasi Sistem (Sulit/Riil)}
\textbf{Skenario:} cari file konfigurasi \texttt{*.conf} pada seluruh sistem,
tetapi buang error permission agar output bersih.

\textbf{Tujuan:} menggunakan \texttt{find} pada root filesystem dan mengelola error dengan redirection.

\begin{lstlisting}[language=bash, caption={Scan *.conf tanpa membanjiri error permission}]
sudo find / -type f -name "*.conf" 2>/dev/null \
  | head -n 30
\end{lstlisting}

\begin{tipbox}
Pola \texttt{2>/dev/null} sangat sering dipakai saat scanning seluruh sistem,
karena banyak direktori memang tidak bisa dibaca user biasa.
\end{tipbox}

\subsection*{Praktikum 4.13 --- locate + updatedb (Mudah)}
\textbf{Tujuan:} menggunakan index database untuk pencarian cepat.

\begin{enumerate}
    \item Pastikan \texttt{plocate} tersedia:
\begin{lstlisting}[language=bash, caption={Memasang locate tool}]
sudo apt update
sudo apt install -y plocate \
  || sudo apt install -y mlocate
\end{lstlisting}

    \item Update database (butuh root), lalu cari:
\begin{lstlisting}[language=bash, caption={Update database locate dan melakukan pencarian}]
sudo updatedb
locate sshd_config | head
\end{lstlisting}
\end{enumerate}

\begin{warningbox}
\texttt{locate} tidak selalu real-time. Jika file baru dibuat dan belum terindeks,
hasil bisa kosong sampai \texttt{updatedb} dijalankan.
\end{warningbox}

\subsection*{Praktikum 4.14 --- Audit Permission Bermasalah dengan find (Kondisi Riil)}
\textbf{Skenario:} admin ingin menemukan file/folder yang ``terlalu terbuka'' (world-writable),
atau file setuid/setgid untuk audit keamanan ringan.

\textbf{Tujuan:} mencari objek yang berisiko dari sisi permission.

\begin{enumerate}
    \item Cari direktori world-writable di home:
\begin{lstlisting}[language=bash, caption={Mencari direktori yang world-writable}]
find ~ -type d -perm -0002 2>/dev/null | head -n 30
\end{lstlisting}

    \item Cari file world-writable:
\begin{lstlisting}[language=bash, caption={Mencari file yang world-writable}]
find ~ -type f -perm -0002 2>/dev/null | head -n 30
\end{lstlisting}

    \item (Opsional) Cari file setuid/setgid:
\begin{lstlisting}[language=bash, caption={Audit setuid/setgid di sistem (butuh sudo)}]
sudo find / -type f -perm -4000 2>/dev/null \
  | head -n 30   # setuid
sudo find / -type f -perm -2000 2>/dev/null \ 
  | head -n 30   # setgid
\end{lstlisting}
\end{enumerate}

\begin{warningbox}
Jangan asal mengubah permission file setuid/setgid sistem. Praktikum ini hanya untuk identifikasi.
\end{warningbox}

% ------------------------------------------------------------------------
\section{Kompresi dan Pengarsipan File (tar, gzip, zip)}
\label{sec:archive}

\subsection*{Penjelasan Konsep}
\textbf{Arsip} (\textit{archive}) menggabungkan banyak file/folder menjadi satu paket.
\textbf{Kompresi} mengecilkan ukuran data.
Pada server, kombinasi umum adalah:
\begin{itemize}
    \item \texttt{tar} untuk mengarsipkan folder
    \item \texttt{gzip} (atau \texttt{bzip2}/\texttt{xz}) untuk mengompresi arsip
    \item \texttt{zip} untuk pertukaran lintas OS (Windows/macOS) atau kebutuhan tertentu
\end{itemize}

\subsection*{tar: flag yang paling sering dipakai}
\begin{itemize}
    \item \texttt{-c} : create (buat arsip)
    \item \texttt{-x} : extract (ekstrak arsip)
    \item \texttt{-t} : list (lihat isi arsip tanpa ekstrak)
    \item \texttt{-f <file>} : nama file arsip (wajib jika pakai file)
    \item \texttt{-v} : verbose (tampilkan file yang diproses)
    \item \texttt{-C <dir>} : ekstrak/kerjakan dari direktori tertentu
    \item \texttt{-z} : gunakan gzip (\texttt{.tar.gz / .tgz})
    \item \texttt{-J} : gunakan xz (\texttt{.tar.xz}) — kompresi lebih kuat, lebih lambat
\end{itemize}

Contoh cepat:
\begin{itemize}
    \item Buat: \texttt{tar -czvf backup.tar.gz folder/}
    \item Lihat isi: \texttt{tar -tzvf backup.tar.gz}
    \item Ekstrak: \texttt{tar -xzvf backup.tar.gz -C /tujuan}
\end{itemize}

\subsection*{gzip: flag penting}
\begin{itemize}
    \item \texttt{-k} : keep file asli (jangan hapus sumber)
    \item \texttt{-d} : decompress (gunzip)
    \item \texttt{-1} sampai \texttt{-9} : level kompresi (1 cepat, 9 maksimum)
\end{itemize}
Catatan: \texttt{gzip} bekerja pada \textbf{file tunggal}. Untuk folder, biasanya \texttt{tar} dulu baru \texttt{gzip}.

\subsection*{zip/unzip: flag penting}
\begin{itemize}
    \item \texttt{zip -r} : rekursif untuk folder
    \item \texttt{zip -9} : kompresi maksimum (lebih lambat)
    \item \texttt{unzip -l} : lihat isi zip tanpa ekstrak
    \item \texttt{unzip -d <dir>} : ekstrak ke direktori tertentu
\end{itemize}

\subsection*{Praktikum 4.15 --- Arsip Workspace: tar + gzip (Mudah)}
\textbf{Tujuan:} membuat \texttt{.tar.gz}, mengecek isi, lalu mengekstraknya.

\begin{enumerate}
    \item Buat arsip terkompresi dari workspace:
\begin{lstlisting}[language=bash, caption={Membuat tar.gz}]
cd ~/praktikum-os
tar -czvf week04-backup.tar.gz week04
ls -lah week04-backup.tar.gz
\end{lstlisting}

    \item Lihat isi arsip tanpa ekstrak:
\begin{lstlisting}[language=bash, caption={Melihat isi tar.gz}]
tar -tzvf week04-backup.tar.gz | head -n 30
\end{lstlisting}

    \item Ekstrak ke folder uji:
\begin{lstlisting}[language=bash, caption={Mengekstrak tar.gz ke folder tertentu}]
mkdir -p ~/praktikum-os/restore-test
tar -xzvf week04-backup.tar.gz \ 
  -C ~/praktikum-os/restore-test
ls -lah ~/praktikum-os/restore-test/week04 | head
\end{lstlisting}
\end{enumerate}

\subsection*{Praktikum 4.16 --- Backup Konfigurasi Sistem (Subset) + Validasi (Kondisi Riil)}
\textbf{Skenario:} sebelum perubahan konfigurasi server, admin membuat backup subset \texttt{/etc} dan menaruhnya di folder backup.

\textbf{Tujuan:} membuat arsip yang rapi, bisa diverifikasi, dan tidak membanjiri error permission.

\begin{enumerate}
    \item Siapkan folder backup:
\begin{lstlisting}[language=bash, caption={Membuat folder backup}]
mkdir -p ~/backup
\end{lstlisting}

    \item Buat backup konfigurasi penting dan verifikasi ukuran:
\begin{lstlisting}[language=bash, caption={Backup konfigurasi penting dengan tar.gz}]
sudo tar -czvf ~/backup/etc-core-$(date +%F).tar.gz \
  /etc/ssh /etc/hosts /etc/netplan /etc/systemd \
  2>/dev/null
ls -lah ~/backup/etc-core-*.tar.gz
\end{lstlisting}

    \item Validasi isi backup (tanpa ekstrak):
\begin{lstlisting}[language=bash, caption={Validasi isi backup}]
tar -tzvf ~/backup/etc-core-$(date +%F).tar.gz \
  | head -n 40
\end{lstlisting}

    \item Restore uji ke direktori aman (jangan overwrite sistem):
\begin{lstlisting}[language=bash, caption={Restore uji ke folder sandbox}]
mkdir -p ~/backup/restore-sandbox
tar -xzvf ~/backup/etc-core-$(date +%F).tar.gz \
  -C ~/backup/restore-sandbox
ls -lah ~/backup/restore-sandbox/etc | head -n 20
\end{lstlisting}
\end{enumerate}

\begin{tipbox}
Pola ``backup subset penting'' lebih aman daripada \texttt{tar /etc} seluruhnya, dan lebih realistis untuk admin.
\end{tipbox}

\subsection*{Praktikum 4.17 --- gzip untuk File Tunggal (Mudah)}
\textbf{Tujuan:} memahami gzip sebagai kompresor file tunggal, serta efek \texttt{-k} dan level kompresi.

\begin{enumerate}
    \item Buat file besar sederhana untuk uji:
\begin{lstlisting}[language=bash, caption={Membuat file uji besar}]
cd ~/praktikum-os/week04
yes "logline" | head -n 200000 > biglog.txt
ls -lah biglog.txt
\end{lstlisting}

    \item Kompres dengan gzip dan simpan file asli:
\begin{lstlisting}[language=bash, caption={gzip dengan keep}]
gzip -k -9 biglog.txt
ls -lah biglog.txt biglog.txt.gz
\end{lstlisting}

    \item Decompress:
\begin{lstlisting}[language=bash, caption={Decompress gzip}]
gzip -d -k biglog.txt.gz
ls -lah biglog.txt biglog.txt.gz
\end{lstlisting}
\end{enumerate}

\subsection*{Praktikum 4.18 --- zip untuk Distribusi ke Lintas OS (Kondisi Riil)}
\textbf{Skenario:} Anda perlu mengirim folder backup/restore ke user Windows/macOS.

\begin{enumerate}
    \item Pastikan zip/unzip ada:
\begin{lstlisting}[language=bash, caption={Install zip/unzip}]
sudo apt update
sudo apt install -y zip unzip
\end{lstlisting}

    \item Buat zip rekursif dan cek isinya:
\begin{lstlisting}[language=bash, caption={Membuat zip dan melihat isi tanpa ekstrak}]
cd ~/backup
zip -r -9 restore-sandbox.zip restore-sandbox
unzip -l restore-sandbox.zip | head -n 30
\end{lstlisting}

    \item Ekstrak ke folder baru:
\begin{lstlisting}[language=bash, caption={Ekstrak zip ke folder tujuan}]
mkdir -p ~/backup/unzip-test
unzip restore-sandbox.zip -d ~/backup/unzip-test
ls -lah ~/backup/unzip-test | head
\end{lstlisting}
\end{enumerate}

% ------------------------------------------------------------------------
\section{Rangkuman}
\label{sec:summary-filesystem}

\begin{itemize}
    \item Struktur Linux berpusat pada \texttt{/} dan mengikuti praktik FHS: konfigurasi di \texttt{/etc}, log di \texttt{/var/log}, home di \texttt{/home}.
    \item Linux punya banyak jenis file: reguler, direktori, symlink, serta file perangkat (\texttt{/dev}) dan virtual FS (\texttt{/proc}, \texttt{/sys}).
    \item Navigasi terminal (\texttt{pwd}, \texttt{ls}, \texttt{cd}, \texttt{pushd/popd}) adalah fondasi administrasi server.
    \item Kontrol akses dibangun dari permission (\texttt{chmod}) dan ownership (\texttt{chown}/\texttt{chgrp}); shared folder adalah skenario nyata yang sering ditemui.
    \item Hard link berbagi inode; symlink menunjuk path target dan cocok untuk pola deploy/rollback.
    \item \texttt{find} fleksibel untuk pencarian real-time; \texttt{locate} cepat tapi tergantung database.
    \item Backup sederhana sangat umum dengan \texttt{tar.gz}; \texttt{gzip} untuk file tunggal; \texttt{zip} nyaman untuk pertukaran lintas OS.
\end{itemize}

% ------------------------------------------------------------------------
\section{Latihan dan Tugas Mandiri}
\label{sec:exercises-filesystem}

\subsection*{A. Latihan Konseptual}
\begin{exercisebox}[4.A]
Pilih 8 direktori dari daftar FHS (\texttt{/etc, /var, /home, /tmp, /usr, /opt, /srv, /dev, /proc, /sys})
dan jelaskan fungsi masing-masing dalam 1--2 kalimat.
\end{exercisebox}

\begin{exercisebox}[4.B]
Jelaskan perbedaan hard link vs symlink, lalu sebutkan 1 use case realistis untuk masing-masing.
\end{exercisebox}

\subsection*{B. Latihan Praktikum (Mudah--Menengah)}
\begin{exercisebox}[4.C]
Buat folder \texttt{\~/praktikum-os/week04-latihan} berisi:
\begin{itemize}
    \item 3 file teks, 1 subfolder, dan 1 symlink yang menunjuk salah satu file.
\end{itemize}
Tampilkan bukti dengan \texttt{ls -lah} dan \texttt{ls -li}.
\end{exercisebox}

\begin{exercisebox}[4.D]
Atur permission file \texttt{laporan.txt} agar:
owner boleh baca+tulis, group hanya baca, others tidak punya akses.
Tuliskan perintahnya dan verifikasi dengan \texttt{ls -l}.
\end{exercisebox}

\begin{exercisebox}[4.E]
Gunakan \texttt{find} untuk mencari 20 file terbesar di \texttt{/var/log}.
Jika permission mengganggu, buang error dengan \texttt{2>/dev/null}.
Tampilkan output 20 baris pertama.
\end{exercisebox}

\subsection*{C. Latihan Praktikum (Kondisi Riil)}
\begin{exercisebox}[4.F --- Kasus Admin 1]
Anda diminta membuat folder \texttt{/srv/data-share} untuk tim (group \texttt{project}) dengan aturan:
\begin{itemize}
\item semua anggota group bisa membuat dan mengedit file
\item file baru mewarisi group folder (setgid)
\item user lain tidak bisa akses
\item user tidak boleh menghapus file milik user lain di folder tersebut (sticky bit)
\end{itemize}
Tuliskan perintah dan bukti \texttt{ls -ld}.
\end{exercisebox}

\begin{exercisebox}[4.G --- Kasus Admin 2]
Anda diminta mengumpulkan bukti troubleshooting:
\begin{itemize}
\item salin 100 baris terakhir \texttt{/var/log/auth.log} ke file di home Anda
\item kompres file tersebut menjadi \texttt{.gz} tanpa menghapus file asli
\item arsipkan folder bukti menjadi \texttt{tar.gz}
\end{itemize}
Tuliskan semua perintah dan tampilkan \texttt{ls -lah} hasil akhirnya.
\end{exercisebox}

\subsection*{Referensi (Manual Pages)}
\begin{itemize}
    \item \texttt{man hier}, \texttt{man ls}, \texttt{man stat}, \texttt{man file}
    \item \texttt{man chmod}, \texttt{man chown}, \texttt{man chgrp}, \texttt{man umask}
    \item \texttt{man ln}
    \item \texttt{man find}, \texttt{man locate}, \texttt{man updatedb}
    \item \texttt{man tar}, \texttt{man gzip}, \texttt{man zip}, \texttt{man unzip}
\end{itemize}
